\section{One coupled field state and one physical energy}
\label{sec:r14-coupling}
The physical unification supplied here is an explicit composition of
electromagnetic, polarization, mechanical and thermal state. Its coupling
terms come from one energy function and its dissipative channels have named
destinations. Maxwell's equations, material constitutive assumptions and
thermodynamic balances keep their established meanings. No additional force
law is inferred from a shared notation.

R13 already supplies separate field and mechanical stores, a Lorentz material
law and a thermoelastic free energy. The missing connection addressed here
is a common evolution when those stores depend on the same material or
mechanical state. Independently changing optical coefficients, adding a force
and depositing absorbed power as heat need not conserve their combined energy.
The following construction makes those dependencies part of the formalization.

\subsection{The common continuous model}
Let $z$ be a finite tuple of nonthermal state coordinates and let
$s=(s_1,\ldots,s_n)$ contain local-equilibrium entropy coordinates, in
$\mathrm{J/K}$. Choose a twice continuously differentiable total energy
$\mathcal H(z,s)$, in joules, and define
\begin{equation}
 e_z=\partial_z\mathcal H,\qquad
 T_i=\partial_{s_i}\mathcal H>0,\qquad
 S_{\rm tot}=\sum_i s_i.
 \label{eq:r14-conjugates}
\end{equation}
These derivatives hold all other state coordinates fixed. Distinct physical
units are retained componentwise; a matrix below represents the corresponding
typed contractions. The temperature is the derivative of the \emph{complete}
energy, including its constitutive coupling, rather than an unrelated lookup.
The use of thermodynamic potentials and their natural-variable derivatives is
established in
\href{https://www.nist.gov/publications/use-legendre-transforms-chemical-thermodynamics}
{UC1: Alberty et al., \emph{Use of Legendre Transforms in Chemical Thermodynamics}}.

Choose $\mathcal J^T=-\mathcal J$, dissipative maps
$\mathcal R_\ell=\mathcal R_\ell^T\succeq0$, an input map $\mathcal G$ and
nonnegative allocation fractions $\alpha_{i\ell}$ with
$\sum_i\alpha_{i\ell}=1$. Let $k_{ij}=k_{ji}\ge0$ be thermal conductances
in $\mathrm{W/K}$, with $k_{ii}=0$, and $Q_i$ be externally supplied heat
rates in watts. Define the versioned coupled model
\begin{align}
 \dot z&=\mathcal J e_z-\sum_\ell\mathcal R_\ell e_z+\mathcal G u,
 &y&=\mathcal G^T e_z,\nonumber\\
 D_\ell&=e_z^T\mathcal R_\ell e_z\ge0,
 &\dot s_i&=\frac{\displaystyle\sum_\ell\alpha_{i\ell}D_\ell
              +\displaystyle\sum_j k_{ij}(T_j-T_i)+Q_i}{T_i}.
 \label{eq:r14-coupled-state}
\end{align}
The input/output pairing $u^Ty$ is supplied nonthermal power. Each $D_\ell$
is power removed from the nonthermal coordinates and thermalized in the
declared stores on the modeled time scale. A channel that stores excitation,
emits radiation or exports power must retain that state or output instead of
being wholly allocated as heat. The heat entropy term $Q_i/T_i$ uses the
receiving node's local-equilibrium temperature. An external reservoir at a
different temperature needs its own entropy flux and interface production;
coherent incident optical power is not assigned a heat entropy $P/T_i$.

This is a selected finite thermodynamic composition. Energy-preserving,
entropy-producing coupling is part of established irreversible
port-Hamiltonian modeling; see
\href{https://arxiv.org/html/2306.08914v2#S2.SS1}
{UC2: Philipp et al., section 2.1}.
The explicit identities below follow directly from the equations here; no
control, stability or convergence result from that source is imported.

\Needspace{15\baselineskip}
\begin{proposition}[Combined first and second laws]
For a differentiable solution in the stated positive-temperature domain,
\begin{align}
 \dot{\mathcal H}&=u^Ty+\sum_i Q_i,\label{eq:r14-total-energy}\\
 \dot S_{\rm tot}-\sum_i\frac{Q_i}{T_i}
 &=\sum_{i,\ell}\frac{\alpha_{i\ell}D_\ell}{T_i}
   +\sum_{i<j}\frac{k_{ij}(T_i-T_j)^2}{T_iT_j}\ge0.
 \label{eq:r14-total-entropy}
\end{align}
\end{proposition}
\begin{proof}
The chain rule gives
$\dot{\mathcal H}=e_z^T\dot z+\sum_i T_i\dot s_i$.
The skew contraction vanishes. Each $-D_\ell$ is canceled by
$\sum_i\alpha_{i\ell}D_\ell$; the two heat rates on every reciprocal
conduction link cancel. This proves the energy identity. Summing the entropy
equations pairs a link's terms as
$k_{ij}[(T_j-T_i)/T_i+(T_i-T_j)/T_j]
=k_{ij}(T_i-T_j)^2/(T_iT_j)$. All remaining production terms are nonnegative.
\end{proof}
Thus an isolated selected model conserves its total energy while dissipation
increases entropy. Energy has not vanished when mechanical or field amplitude
decays: it has entered the same model's thermal state. Positivity hypotheses
are constitutive assumptions, not conclusions about every proposed device.

The maps may depend smoothly on state. Local Lipschitz regularity of the
resulting right-hand side, regular finite coordinates, admissible initial data
and a prescribed regular input give the usual local ODE obligation. These
conditions do not prove global continuation, material stability, or that a
trajectory remains in the positive-temperature domain. Lower-bounded energy,
constitutive stability and any further operating-region invariance are checked
for the selected application. Explicit time dependence adds
$\partial_t\mathcal H$ to equation~\eqref{eq:r14-total-energy}.

\subsection{A concrete Maxwell--Lorentz, mechanical and thermal specialization}
Select a finite electromagnetic representation on a fixed reference topology
with power-dual coordinates. Let $d$ be electric displacement-flux coordinates
in coulombs, $b$ magnetic flux coordinates in webers, and $\mathsf C$
a fixed signed curl/interconnection matrix. Their conjugates $e,h$ have
voltage and current units. Matrix entries and coordinate normalization are
part of the selected field reduction; these are not point samples or pixel
intensities. The skew pair $(\mathsf C,-\mathsf C^T)$ is the finite
counterpart of curl integration by parts, with boundary power supplied through
explicit ports. A compatible field reduction must derive this pair and its
boundary terms. Such preservation of Maxwell power structure is developed in
\href{https://arxiv.org/html/2202.04390v2}
{UC3: Brugnoli, Rashad and Stramigioli, \emph{Dual field structure-preserving
discretization of port-Hamiltonian systems}}.

Add mechanical coordinates $q$, momenta $p$, polarization-charge coordinates
$P_j$ of the same dimension as $d$, and their conjugate flux coordinates
$\pi_j$. Here $q$ is a local mechanical coordinate, distinct from the digital
JK word $q_n$. Set $\delta=d-\sum_j P_j$ and choose
\begin{align}
 \mathcal H={}&\tfrac12p^TM^{-1}p+U(q,s)
       +\tfrac12\delta^T C_e(q,s)^{-1}\delta
       +\tfrac12 b^TL_b(q,s)^{-1}b\nonumber\\
       &+\sum_{j=1}^{m}\left(
          \tfrac{f_j}{2}\pi_j^T\pi_j
          +\tfrac{\omega_j^2}{2f_j}P_j^TP_j\right).
 \label{eq:r14-concrete-energy}
\end{align}
Here $M=M^T\succ0$ is fixed, $C_e=C_e^T\succ0$ and $L_b=L_b^T\succ0$
are selected electric and magnetic energy matrices, $m$ is finite, and
$f_j>0$, $\omega_j>0$ are fixed oscillator coefficients. Under the stated
flux/charge normalization $C_e,L_b,f_j$ have units
$\mathrm F,\mathrm H,\mathrm{F/s^2}$; $\pi_j$ has units $\mathrm{V\,s}$.
These finite coefficients are projected/material parameters and are not
numerically identified with the continuum coefficients of
equation~\eqref{eq:r13-em-lorentz} without the actual normalization map.
Mixed mechanical dimensions retain the inherited typed-block convention.

The internal energy $U(q,s)$ includes elastic/thermal coupling once. In
particular it may be the entropy-variable form of R13's thermoelastic free
energy, rather than an independently added second elastic store. Define
\begin{align}
 v&=M^{-1}p,& e&=C_e^{-1}\delta,& h&=L_b^{-1}b,\nonumber\\
 V_j&=f_j\pi_j,&
 w_j&=-e+\frac{\omega_j^2}{f_j}P_j,&
 T_i&=\partial_{s_i}\mathcal H.
 \label{eq:r14-concrete-efforts}
\end{align}
For symmetric positive-semidefinite mechanical damping $C_m$ and electric
conductance $G_e$, and $\gamma_j\ge0$, the nonthermal equations are
\begin{align}
 \dot q&=v,&\dot p&=-\partial_q\mathcal H-C_m v+F_{\rm ext},\nonumber\\
 \dot d&=\mathsf C h-G_e e+i_{\rm ext},&
 \dot b&=-\mathsf C^T e,\nonumber\\
 \dot P_j&=V_j,&
 \dot\pi_j&=-w_j-\frac{\gamma_j}{f_j}V_j.
 \label{eq:r14-concrete-evolution}
\end{align}
The displayed supply is $v^TF_{\rm ext}+e^Ti_{\rm ext}$.
Additional selected electromagnetic boundary ports are paired by the same
power convention. The dissipation channels deposited through
equation~\eqref{eq:r14-coupled-state} are exactly
\begin{equation}
 D_m=v^TC_m v,\qquad D_e=e^TG_e e,\qquad
 D_j=\frac{\gamma_j}{f_j}V_j^TV_j.
 \label{eq:r14-concrete-loss}
\end{equation}
The curl powers cancel as $e^T\mathsf C h-h^T\mathsf C^Te=0$;
each mechanical or polarization canonical pair cancels in the same way.
The remaining nonthermal power is the supply minus the displayed losses,
so the preceding energy and entropy theorem applies without an additional
unidentified transfer coefficient.

For the fixed $f_j$, differentiating $V_j=f_j\pi_j$ gives
\begin{equation}
 \dot V_j=f_j e-\omega_j^2P_j-\gamma_jV_j.
 \label{eq:r14-finite-lorentz}
\end{equation}
This is the finite-coordinate Lorentz oscillator. The same $P_j$ enters
$\delta$ in field energy, supplies the electromagnetic response and receives
the work that its damping converts to heat. No separate full ``absorption
heat'' is added on top of $D_j$. The constitutive origin and continuum energy
proof remain those of equations~\eqref{eq:r13-em-lorentz}--\eqref{eq:r13-em-storage}.

\subsection{Reciprocal optical force and thermal response}
For fixed $d,b,P_j,\pi_j,p,s$, matrix inverse differentiation gives the
mechanical load contributed by the two field stores:
\begin{equation}
 (F_{\rm field})_a
  =\tfrac12 e^T(\partial_{q_a}C_e)e
   +\tfrac12 h^T(\partial_{q_a}L_b)h.
 \label{eq:r14-field-force}
\end{equation}
Thus $-\partial_{q_a}\mathcal H=-\partial_{q_a}U+(F_{\rm field})_a$
for the selected fixed $M,f_j,\omega_j$. This derives the reciprocal force
from the same geometry/material dependence that changes the field response.
It does not insert a generic radiation-pressure formula into every medium.
The scalar electroquasistatic restriction is the familiar
$F=(e^2/2)\,\partial_q C_e$, with energy differentiation performed at fixed
charge, as derived in
\href{https://ocw.mit.edu/courses/6-641-electromagnetic-fields-forces-and-motion-spring-2009/044362926d4b6bc9f849f4ece0c4c1f7_MIT6_641s09_lec12.pdf}
{UC4: MIT Electromagnetic Fields, Forces, and Motion, lecture 12, pages 1--2}.
Holding voltage by a source changes the charge and supplies electrical work;
one cannot differentiate a fixed-voltage expression as though that source
were absent.

Likewise the common thermal derivative is
\begin{equation}
 T_i=\partial_{s_i}U
       -\tfrac12 e^T(\partial_{s_i}C_e)e
       -\tfrac12 h^T(\partial_{s_i}L_b)h.
 \label{eq:r14-field-temperature}
\end{equation}
The positivity domain applies to this complete expression. Simply prescribing
both arbitrary $C_e(q,s)$ and an unrelated $T_i(s_i)$ would overdetermine the
thermodynamic constitutive model. As a simple uncoupled thermal store,
$U_i(s_i)=c_iT_{i0}\exp[(s_i-s_{i0})/c_i]$ with $c_i,T_{i0}>0$ gives
$\partial_{s_i}U_i=T_{i0}\exp[(s_i-s_{i0})/c_i]>0$; the exponent is
dimensionless. Once field or elastic storage depends on entropy, its derivative
must also be included. This identity is a constitutive illustration, not a
numerical example or a universal material law.

Material data are often supplied versus temperature. For one thermal
coordinate, a twice differentiable Helmholtz energy $\mathcal F(z,T)$ gives
\begin{equation}
 s=-\partial_T\mathcal F(z,T),\qquad
 c_z=-T\partial_T^2\mathcal F>0,\qquad
 \mathcal H(z,s)=\mathcal F(z,T(z,s))+T(z,s)s.
 \label{eq:r14-legendre}
\end{equation}
At fixed $z$, $\partial_Ts=c_z/T>0$ permits a local inverse on the declared
branch. Differentiation cancels the $\partial T$ terms and yields
$\partial_s\mathcal H=T$ and
$\partial_z\mathcal H=\partial_z\mathcal F|_T$. For several temperatures the
corresponding $-\partial_T^2\mathcal F$ must be nonsingular on the selected
branch; positive definiteness supplies local thermodynamic stability of that
temperature block. A constitutive inversion still needs its supplied exact
finite definition and domain. This connects R13's free-energy closure to the
common entropy state without omitting its reciprocal thermoelastic term.

The force theorem is a theorem about the chosen reduced energy. To represent
a moving physical device, $q$ must describe a declared deformation or component
coordinate with a valid energy reduction and boundary work. Fixed reference
topology with parameterized energy matrices can describe selected small-motion
or quasistatic couplings; it is not an arbitrary moving-medium Maxwell model.
Unmodeled basis motion, motional electromotive terms, relativistic effects or
unrepresented mechanical momentum are not supplied by writing a derivative.

\subsection{Field constraints, material modes and exact seed binding}
The Maxwell divergence conditions still constrain any proposed field
representation. For an unreduced compatible flux description, let
$N_e\mathsf C=0$ and $N_b\mathsf C^T=0$ express the boundary-of-boundary
identities. If $N_ed=\rho$ initially and the free-charge coordinates obey
\begin{equation}
 \dot\rho=N_e(-G_e e+i_{\rm ext}),\qquad N_bb=0\ \hbox{initially},
 \label{eq:r14-charge-compatible}
\end{equation}
then differentiation preserves both constraints. Compatible boundary ports
contribute their corresponding charge term. The signs use
$i_{\rm ext}$ as an injected electric-flow input; an impressed current leaving
the field has the opposite sign. Polarization is already part of $d$ and is
not an additional free-charge source.

For the exact finite ODE binding, the selected representation builds these
constraints into a regular basis or supplies their explicit elimination and
compatible independent inputs, as required by R13. The invariance calculation
does not solve incompatible initial data, provide an unspecified boundary
condition, or add a differential-algebraic solver to XOP. Modal truncation,
constitutive fit and the mapping from finite fluxes to physical fields remain
explicit approximation/evidence obligations.

ASA/NA and JK can select a constitutive mode $\chi$, an actuator input or a
boundary configuration. In a fixed mode the equations use
$\mathcal H(z,s;\chi)$. At a mode change with unchanged physical coordinates,
\begin{equation}
 W_{\rm switch}=\mathcal H(z,s;\chi^+)-\mathcal H(z,s;\chi^-).
 \label{eq:r14-mode-energy}
\end{equation}
The value is a required signed supply or export, not proof that an actuator
can deliver it. A reset of coordinates or entropy has its own complete
before/after balance. A dynamic reporter or hysteretic state affecting energy
must likewise have its conjugate exchange and evolution; a declared empirical
readout may remain an observation without being counted as an energy store.

The seed specialization records $\mathcal H$, its explicit conjugate bodies,
the selected matrices, thermal allocations, inputs, initial state and domains
in one typed operator graph. Existing scalar arithmetic, finite sums,
\code{EXP}, and nonsingular \code{SOLVE\_LINEAR} express the displayed finite
bodies when their operands meet the inherited contracts. Gradients in this
chapter are mathematical definitions; the specialization supplies their
derived finite bodies. Existing \code{DIFF} can denote a derivative on its
declared domain; it does not by itself provide an automatic differentiation
algorithm. The finite ODE notation retains the inherited existence and denotation
contract. One-bit words preserve this authoritative graph and its operands.

R13's midpoint balance applies to its fixed linear mechanical specialization.
It does not automatically extend to the nonlinear coupled energy
\eqref{eq:r14-concrete-energy}. Any selected finite update for the combined
model needs its own discrete work, heat and entropy identities or stated
defects. This chapter adds the integrated physical equations and their
continuous proofs; it does not silently claim a new time-step implementation.
